\section{Lindstr\"om quantifiers and back-and-forth}
\label{sec:p03-lindstrom-quantifiers}

A generalized quantifier tests a structure defined by its argument
formulas. The defining class is part of the quantifier's semantics;
it need not be first-order definable. This description separates three
questions: what class the quantifier recognizes, which closure operations
its logic retains, and which partial isomorphisms preserve its formulas.

\begin{hbconvention}{def:p03-lq-setting}{Structures and formula parameters}
Work with one-sorted structures with nonempty domains, in classical set
theory with Choice. All function and relation arities are finite.
The first-order and \(L_{\infty\omega}\) languages used below have
finitely many free variables in each formula. For abstract logics,
expressive equivalence concerns sentences over the same vocabulary.
The translations also preserve formulas with finitely many free
parameters, represented by fresh constants as in
\cref{def:p03-w03-p03-047-closure}. No claim is made that a formula
with infinitely many free parameters has a finite-parameter translation.
\end{hbconvention}

\subsection{Quantifiers from classes of structures}

\begin{hbdefinition}{def:p03-lq-full}{A full-domain class quantifier}
Let \(\sigma\) be a finite vocabulary, called the \emph{pattern}, and
let \(K\) be an isomorphism-closed class of \(\sigma\)-structures.
For each relation \(R\) of arity \(r\), an application of \(Q_K\)
has an operand \(\alpha_R(\bar x_R,\bar u)\) with \(r\) bound
variables. For each function \(f\) of arity \(r\), it has a graph
operand \(\beta_f(\bar x_f,y_f,\bar u)\), and for each constant
\(c\), a singleton operand \(\gamma_c(z_c,\bar u)\).
Each bound list consists of distinct variables; these lists are
chosen disjoint from the free parameter list \(\bar u\).
Write the application as
\[
 Q_K\bigl((\bar x_R\,\alpha_R)_R,
          (\bar x_f y_f\,\beta_f)_f,(z_c\,\gamma_c)_c\bigr).
\]
Each list binds variables only in its own operand.
In a structure \(\mathfrak A\), with parameters \(\bar a\), it is
true precisely when the operands define, on the whole domain \(A\),
the relations, total function graphs and unique constants of a
\(\sigma\)-structure belonging to \(K\).
If a function operand is not a total single-valued graph, or a
constant operand is not a singleton, the application is false.
A nullary relation has an operand with no bound variables.

For a family of such quantifiers, \(\operatorname{FO}(Q_i:i\in J)\)
is the ordinary finite formula language obtained from first-order
atoms by Boolean operations, first-order quantifiers and these
applications. Each formula uses only finitely many members of the
family. With a general base logic, \(\mathscr L(Q_i:i\in J)\)
is the analogous finite closure starting from its formulas.
This is the class-quantifier construction of
\cite[Ch.~II, Definition~4.1.1, pp.~49--50]{Ebbinghaus1985ExtendedLogics}.
\end{hbdefinition}

\begin{hbexample}{def:p03-lq-examples}{Standard quantifiers as class tests}
The following choices of \(K\) recover familiar operations directly
from \cref{def:p03-lq-full}.
\begin{enumerate}
 \item For a unary predicate \(P\), take the structures with
       \(P\ne\varnothing\). Then \(Q_Kx\,\alpha\) is
       \(\exists x\,\alpha\).
 \item For an \(r\)-ary predicate \(R\), with \(r\geq1\), take the structures
       containing \(C\subseteq A\) with \(|C|\geq\aleph_\alpha\)
       and \(C^r\subseteq R\). The result is the homogeneous-set
       quantifier \(Q^r_\alpha\); at \(r=1\) it is the usual
       cardinality quantifier.
 \item For binary \(R\), take the strict linear orders of the whole
       domain with cofinality \(\omega\). This is
       \(Q^{\mathrm{cf}\,\omega}\). Taking all well-orders of the
       whole domain instead gives \(Q^{\mathrm{wo}}\).
 \item For unary \(P\), take \(|P|=|A|\). This is Chang's
       full-domain quantifier \(Q^{\mathrm C}\).
\end{enumerate}
In the order examples, an operand defining a relation that is not a
linear order is rejected. The last example compares with the whole
ambient domain, not a second supplied predicate
\cite[Ch.~II, Example~4.1.2, p.~50]{Ebbinghaus1985ExtendedLogics}.
\end{hbexample}

\begin{hbtheorem}{thm:p03-lq-representation}{proved=annex:p03.lq.1}{Representation by class quantifiers}
Let \(\mathscr L\) be a one-sorted regular logic.
\begin{enumerate}
 \item Every isomorphism-closed class \(K\) of a finite pattern
       vocabulary is defined by one sentence of
       \(\operatorname{FO}(Q_K)\).
 \item If \(K\) is already defined by a sentence of \(\mathscr L\),
       every formula of \(\mathscr L(Q_K)\) has an equivalent
       \(\mathscr L\)-formula.
 \item Suppose \(\mathscr L\) is \emph{finitary}: each formula
       belongs to the logic of a finite sub-vocabulary, with its
       finite list of parameter constants retained. Then
       \(\mathscr L\) is expressively equivalent to first-order
       logic with the quantifiers \(Q_K\) for all its definable
       classes of finite vocabulary.
\end{enumerate}
The family in the last clause need not be finite
\cite[Ch.~II, Theorem~4.1.3, p.~51]{Ebbinghaus1985ExtendedLogics}.
\end{hbtheorem}

\begin{hbtheorem}{thm:p03-lq-second-order}{proved=annex:p03.lq.2}{Second-order logic needs an infinite family}
Full second-order logic is not expressively equivalent to
\(\operatorname{FO}(Q_1,\ldots,Q_s)\) for any finite family of
ordinary Lindstr\"om quantifiers of finite pattern vocabulary.
Consequently the representation of second-order logic in
\cref{thm:p03-lq-representation} cannot use a finite family.
This is the non-finite-generation assertion following
\cite[Ch.~II, Theorem~4.1.3, p.~51]{Ebbinghaus1985ExtendedLogics}.
\end{hbtheorem}

\subsection{Relativized domains}

\begin{hbdefinition}{def:p03-lq-relative}{A domain-relativized class quantifier}
An application of \(Q_K^*\) has, in addition to the operands of
\cref{def:p03-lq-full}, a domain operand \(\delta(v,\bar u)\).
In \(\mathfrak A\) its domain is
\(D=\{b\in A:\mathfrak A\models\delta(b,\bar a)\}\).
The application is true exactly when \(D\ne\varnothing\) and the
other operands, evaluated in \(\mathfrak A\) but restricted to tuples
from \(D\), define a \(\sigma\)-structure in \(K\).
Function operands must be total single-valued graphs on \(D\), and
constant operands must define singletons in \(D\).
No condition is imposed on those operands outside \(D\).
This is the domain variant of
\cite[Ch.~II, Definition~4.1.4, p.~51]{Ebbinghaus1985ExtendedLogics}.
\end{hbdefinition}

\begin{hbtheorem}{thm:p03-lq-domain}{proved=annex:p03.lq.3}{Domain translation and regularity}
Let \(K\) have finite pattern \(\sigma\), and let \(U\) be a fresh
unary relation. Define \(K^U\) to consist of the
\((\sigma\cup\{U\})\)-structures in which \(U\) is nonempty and
\(\sigma\)-closed and the induced \(\sigma\)-structure on \(U\)
belongs to \(K\). Then:
\begin{enumerate}
 \item \(Q_K\) is expressed by \(Q_K^*\) with domain formula \(v=v\).
 \item The logics \(\operatorname{FO}(Q_K^*)\) and
       \(\operatorname{FO}(Q_{K^U})\) are expressively equivalent.
       The translations preserve free parameters and reject malformed
       function graphs and constant descriptions correctly.
 \item For any family of finite patterns,
       \(\operatorname{FO}(Q_i^*:i\in J)\) is regular, with
       substitution interpreted on admissible interpretations as in
       \cref{def:p03-w03-p03-047-relativization}.
\end{enumerate}
The class conversion is
\cite[Ch.~II, Proposition~4.1.5, p.~52]{Ebbinghaus1985ExtendedLogics}.
\end{hbtheorem}

\begin{hbcorollary}{thm:p03-lq-minimal}{proved=annex:p03.lq.4}{The regular extension generated by a class}
Let \(K\) have finite pattern vocabulary. Every regular logic in
which \(K\) is definable contains the expressive power of
\(\operatorname{FO}(Q_K^*)\). Thus this logic is the least regular
extension of first-order logic defining \(K\), up to expressive
equivalence. More generally, if \(\mathscr L(Q_K^*)\) is itself
regular, it is the least regular extension of \(\mathscr L\)
defining \(K\).
\end{hbcorollary}

\begin{hbtheorem}{thm:p03-lq-chang-boundary}{proved=annex:p03.lq.5}{A failure of relativization}
The logic \(\operatorname{FO}(Q^{\mathrm C})\) is not closed under
relativization to a nonempty unary domain. There are structures in the
vocabulary \(\{U,V\}\) satisfying the same sentences of this logic,
but differing on whether \(|U\cap V|=|U|\).
Thus ordinary full-domain class quantifiers do not automatically give
regular extensions, even when their defining classes are monotone.
\end{hbtheorem}

\subsection{A standard part is not a quantifier}

\begin{hbdefinition}{def:p03-lq-standard-part}{Standard-part semantics}
Fix a finite pattern \(\sigma\), a fresh unary \(U\), and an
isomorphism-closed class \(K\) of \(\sigma\)-structures.
The standard-part semantics \(\operatorname{FO}[K]\) uses ordinary
first-order sentences over vocabularies containing \(\sigma\cup\{U\}\),
but permits only structures in which \(U\) is nonempty and
\(\sigma\)-closed and its induced \(\sigma\)-structure belongs to \(K\).
On permitted structures satisfaction is ordinary first-order satisfaction.
This is the one-sorted standard-part construction of
\cite[Ch.~II, \S2.6, pp.~39--40]{Ebbinghaus1985ExtendedLogics}.
The brackets distinguish a restriction on permitted structures from
the operator extension \(\operatorname{FO}(Q_K^*)\).
\end{hbdefinition}

\begin{hbtheorem}{thm:p03-lq-standard-omega}{proved=annex:p03.lq.6}{Standard parts and projective representation}
\begin{enumerate}
 \item Every sentence of \(\operatorname{FO}[K]\), with its
       standard-part satisfaction condition, has an equivalent
       sentence of \(\operatorname{FO}(Q_K^*)\).
 \item Let \(K_\omega\) be the class of orders isomorphic to
       \((\omega,<)\). For a vocabulary disjoint from the fixed
       standard-part symbols, the infinite models of any
       \(\operatorname{FO}(Q_{K_\omega}^*)\)-sentence form a PC
       class in \(\operatorname{FO}[K_\omega]\).
       Its models of all nonempty cardinalities form an RPC class
       in that standard-part semantics.
 \item Under the same-domain definition of PC in
       \cref{def:p03-w03-p03-047-ec-pc-rpc}, the second clause cannot
       be strengthened to a PC representation of all nonempty models:
       even a tautology's finite models cannot be such reducts.
\end{enumerate}
The finite-domain qualification is essential for the one-sorted
projective comparison discussed in
\cite[Ch.~II, \S4.1, p.~52]{Ebbinghaus1985ExtendedLogics}.
\end{hbtheorem}

\begin{hbtheorem}{thm:p03-lq-field-boundary}{proved=annex:p03.lq.7}{Different compactness behaviours}
Let \(K\) be the class of characteristic-zero fields in the usual
finite field vocabulary. Then \(\operatorname{FO}[K]\) is compact,
whereas \(\operatorname{FO}(Q_K^*)\) is not countably compact.
The failure in the latter has witnesses to finite satisfiability
with infinite domains. Thus restricting models to have a standard
part and adding the corresponding class quantifier are different
operations even at the level of compactness
\cite[Ch.~II, \S4.1, p.~52]{Ebbinghaus1985ExtendedLogics}.
\end{hbtheorem}

\subsection{Monotone quantifiers and finite comparisons}

\begin{hbdefinition}{def:p03-lq-monotone}{Monotone classes and partial maps}
For a finite relational pattern \(\sigma\), a class \(K\) is
\emph{monotone} if enlarging any or all of its relations on a fixed
domain preserves membership in \(K\). The relation interpretations
are enlarged componentwise; the domain is not changed.

A \emph{graph partial isomorphism} \(p\) from \(\mathfrak A\) to
\(\mathfrak B\) is an injective map from a subset of \(A\) into \(B\)
which preserves and reflects the following atoms whenever all their
displayed variable values lie in its domain:
\[
 R(\bar x),\qquad x=y,\qquad c=y,\qquad f(\bar x)=y.
\]
Nullary relation atoms are included. The domain need not be closed
under functions or contain the constants. A formula is
\emph{term-reduced} if these are its only atomic forms.
Its quantifier rank is zero at atoms, unchanged by negation, the
maximum at a finite Boolean combination, and increased by one at
an existential quantifier or a class-quantifier application.
At the latter, use the maximum of the ranks of its operands;
the maximum of an empty list is zero. Empty conjunction and disjunction
are the rank-zero logical constants true and false.
\end{hbdefinition}

\begin{hbdefinition}{def:p03-lq-systems}{Finite and unbounded comparison systems}
Fix a family of monotone quantifiers with finite relational patterns.
An \(n\)-system between \(\mathfrak A\) and \(\mathfrak B\) is a
sequence \(I_0,\ldots,I_n\) of nonempty sets of finite graph partial
isomorphisms, satisfying the following for every \(p\in I_{m+1}\),
where \(m<n\).
\begin{enumerate}
 \item For each \(a\in A\) there is \(q\in I_m\) extending \(p\)
       with \(a\) in its domain; symmetrically every \(b\in B\)
       can be included in the range of such an extension.
 \item For each quantifier \(Q_K\) and each pattern structure
       \(\mathfrak C\in K\) with domain \(A\), there is
       \(\mathfrak D\in K\) with domain \(B\) such that, for every
       pattern relation \(R\) and \(\bar d\in R^{\mathfrak D}\),
       some \(q\in I_m\) extends \(p\), includes all entries of
       \(\bar d\) in its range, and has
       \(q^{-1}(\bar d)\in R^{\mathfrak C}\).
       The symmetric condition, starting with \(\mathfrak D\) on
       \(B\), is also required.
\end{enumerate}
An \(\omega\)-system is a sequence \((I_m)_{m<\omega}\) satisfying
these conditions at every finite level. An \emph{unbounded system}
is a nonempty set \(I\) of finite graph partial isomorphisms with
the same extension conditions, but with every responding map again
in \(I\), without decreasing a level.
The empty tuple in a true nullary relation is treated by the same
rule; no tuple entries then have to be added.
\end{hbdefinition}

\begin{hbtheorem}{thm:p03-lq-finite-comparison}{proved=annex:p03.lq.8}{Finite comparison theorem}
Let the ambient vocabulary be finite and let \(Q_1,\ldots,Q_s\)
be a finite family of monotone quantifiers with finite relational
patterns. For every \(n<\omega\), the following are equivalent:
\begin{enumerate}
 \item \(\mathfrak A\) and \(\mathfrak B\) agree on all term-reduced
       sentences of \(\operatorname{FO}(Q_1,\ldots,Q_s)\) of rank
       at most \(n\).
 \item There is an \(n\)-system between them.
\end{enumerate}
Consequently they agree on all sentences of this logic if and only if
there is an \(n\)-system for every finite \(n\), and if and only if
there is an \(\omega\)-system. Arbitrary terms are allowed in this
last sentence-equivalence assertion: term reduction preserves truth,
but can increase rank.
For no added quantifiers this recovers the first-order comparison
theorem. This is the monotone-quantifier comparison of
\cite[Ch.~II, Theorem~4.2.3 and Corollary~4.2.4,
pp.~54--56]{Ebbinghaus1985ExtendedLogics}.
\end{hbtheorem}

\begin{hbdefinition}{def:p03-lq-characteristic}{The source's characteristic formulas}
Fix a finite ambient vocabulary and one monotone quantifier \(Q_K\)
for a single relation of arity \(r\geq1\). For \(\bar a\in A^k\),
let \(\xi^0_{\bar a}(\bar x)\) be its complete graph-atomic diagram
in the displayed variables: the conjunction of true graph atoms and
negations of false ones. For a set \(M\) of tuples of one fixed
finite length, abbreviate
\[
 O^m_{\bar a,M}(\bar x,\bar y)
       =\bigvee_{\bar c\in M}\xi^m_{\bar a\bar c}(\bar x,\bar y).
\]
For length one, write \(A\) for its set of one-tuples. Define
\[
\begin{aligned}
 \xi^{m+1}_{\bar a}(\bar x)={}&
 \bigwedge_{c\in A}\exists y\,\xi^m_{\bar a c}(\bar x,y)
       \\
 &\land\forall y\,O^m_{\bar a,A}(\bar x,y)\\
 &\land\!\bigwedge_{\substack{M\subseteq A^r\\ (A,M)\in K}}
             Q_K\bar y\,O^m_{\bar a,M}(\bar x,\bar y)\\
 &\land\!\bigwedge_{\substack{M\subseteq A^r\\ (A,M)\notin K}}
             \neg Q_K\bar y\,\neg
                     O^m_{\bar a,A^r\setminus M}(\bar x,\bar y).
\end{aligned}
\]
All displayed indexed conjunctions and disjunctions are replaced by
finite lists of their equivalence-distinct operands. The finite-type
argument in \cref{annex:p03.lq.8} justifies these finite choices;
no infinitary connective is being added. Retain the displayed
quantifier wrappers, even when a formula happens to be a tautology.
This is the construction denoted \(\psi^m\) in
\cite[Ch.~II, Definition~4.2.5, p.~55]{Ebbinghaus1985ExtendedLogics}.
\end{hbdefinition}

\begin{hbtheorem}{thm:p03-lq-characteristic}{proved=annex:p03.lq.14}{Characteristic formulas and types}
The formula \(\xi^m_{\bar a}\) has rank \(m\), is true at
\(\bar a\) in \(\mathfrak A\), and defines the complete
rank-at-most-\(m\) term-reduced type of that tuple.
Consequently
\(\xi^m_{\bar a\bar c}(\bar x,\bar y)\models\xi^m_{\bar a}(\bar x)\)
and
\(\xi^{m+1}_{\bar a}\models\xi^m_{\bar a}\).
Thus these formulas are equivalent to the characteristic formulas
used in \cref{annex:p03.lq.8}, with the source's syntactic rank and
projection properties retained
\cite[Ch.~II, Lemma~4.2.6, p.~56]{Ebbinghaus1985ExtendedLogics}.
\end{hbtheorem}

\subsection{Infinitary comparisons}

\begin{hbdefinition}{def:p03-lq-infinitary}{Infinitary formulas with class quantifiers}
The language \(L_{\infty\omega}(Q_i:i\in J)\) additionally admits
conjunctions and disjunctions indexed by arbitrary sets, provided the
union of their free-variable sets is finite. Quantifier applications
retain finite pattern vocabularies and finite bound-variable lists.
Formula trees are well-founded. The ambient vocabulary and the
family of quantifiers need not be finite.
\end{hbdefinition}

\begin{hbtheorem}{thm:p03-lq-karp}{proved=annex:p03.lq.9}{Infinitary comparison theorem}
For any vocabulary and any family of monotone finite relational
class quantifiers, two structures satisfy the same sentences of
\(L_{\infty\omega}(Q_i:i\in J)\) if and only if there is an
unbounded comparison system between them. This is the generalized
Karp theorem of \cite[Ch.~II, Theorem~4.3.3,
p.~59]{Ebbinghaus1985ExtendedLogics}.
\end{hbtheorem}

\begin{hbcorollary}{thm:p03-lq-countable}{proved=annex:p03.lq.10}{Ordinary comparison systems}
With no added quantifiers, an unbounded graph comparison system is
an ordinary partial isomorphism system in the sense of
\cref{def:p03-w03-p03-039-partial-iso}. In particular, two at most
countable structures linked by such a system are isomorphic.
This countability condition cannot be omitted.
Moreover, if a logic \(\mathscr R\) gives the same sentences to any
two ordinarily partially isomorphic structures, then
\(\mathscr R\leq_{\equiv}L_{\infty\omega}\) in the sense of
\cref{def:p03-w03-p03-047-strength}. This is a comparison of
which structures can be distinguished, not a claimed translation
of every \(\mathscr R\)-sentence into one infinitary sentence.
\end{hbcorollary}

\subsection{Uncountability moves and interpolation}

\begin{hbdefinition}{def:p03-lq-qone-game}{The finite uncountability game}
The game for \(Q_1\) has \(n\) rounds and a current finite map of
selected elements. In each round the first player chooses either:
\begin{enumerate}
 \item an element in either structure; the second player answers
       with an element in the other; or
 \item an uncountable subset of either structure; the second player
       answers with an uncountable subset of the other, the first
       player chooses an element of this answering subset, and the
       second player chooses its mate from the original subset.
\end{enumerate}
The selected pair extends the current map. The second player wins
if the resulting map after all rounds is a graph partial isomorphism;
an inconsistent or noninjective choice loses immediately.
If no answering uncountable subset exists, the second player loses
that move. This is the game described in
\cite[Ch.~II, \S4.2, p.~57]{Ebbinghaus1985ExtendedLogics}.
\end{hbdefinition}

\begin{hbtheorem}{thm:p03-lq-qone-game}{proved=annex:p03.lq.11}{Games and uncountability}
In finite vocabulary, two structures agree on all sentences of
\(\operatorname{FO}(Q_1)\) if and only if the second player has a
winning strategy in the finite uncountability game for every
\(n<\omega\).
\end{hbtheorem}

\begin{hbtheorem}{thm:p03-lq-qone-interpolation}{proved=annex:p03.lq.12}{Failure of interpolation for uncountability}
Let \(\mathfrak A\) and \(\mathfrak B\) be equivalence-relation
structures with infinitely many classes, each class uncountable.
There is an unbounded \(Q_1\)-comparison system between them.
Consequently they agree on every sentence of
\(L_{\infty\omega}(Q_1)\), regardless of the number of their classes.
In particular, the entailment of
\cref{def:p03-w03-p03-047-keisler-entailment} has no interpolant
in \(\operatorname{FO}(Q_1)\).
This supplies the comparison underlying
\cite[Ch.~II, Application~4.2.7, p.~57]{Ebbinghaus1985ExtendedLogics}.
\end{hbtheorem}

\subsection{Ordinary comparison systems in algebra and order}

\begin{hbtheorem}{thm:p03-lq-examples}{proved=annex:p03.lq.13}{Examples of ordinary comparison systems}
\begin{enumerate}
 \item Any two infinite structures of the empty vocabulary have an
       unbounded ordinary partial isomorphism system.
 \item Any two dense linear orders without endpoints have such a
       system. For every uncountable cardinal \(\kappa\), there
       are two nonisomorphic such orders of cardinality \(\kappa\).
 \item Any two algebraically closed fields of the same characteristic,
       each of infinite transcendence degree over its prime field,
       have such a system.
\end{enumerate}
These are the ordinary partial-isomorphism examples in
\cite[Ch.~II, \S4.3, pp.~57--58]{Ebbinghaus1985ExtendedLogics}.
\end{hbtheorem}

